% !PRINTED 441
\weberpage{441}
%<<CONT
$\dfrac{n^2-1}{8}$ is an even number, in the last two an odd one,
and accordingly we obtain the exact formula (1):

\weq{3}{2^{\frac{n-1}{2}} \overset{\nu}{\Pi} \cos \frac{2\nu m \pi}{n} = (-1)^{\frac{n^2-1}{8}}.}

In formula (2) the sign depends on $m$; it is positive or
negative according as, among the numbers $(\varrho)$, an even or an
odd number of negative ones occurs.

Now, in order to express in the notation the dependence of
the sign on $m$ and $n$, we write for (2)

\weq{4}{2^{\frac{n-1}{2}} \overset{\nu}{\underset{1,\,\frac{n-1}{2}}{\Pi}} \sin \frac{2\nu m \pi}{n} = \left(\frac{m}{n}\right)\sqrt{n},}

in which

\[
\left(\frac{m}{n}\right) = \pm 1
\]

holds.

The symbol $\left(\dfrac{m}{n}\right)$ has been introduced into number theory
from another side, and in a more special form, by \gesperrt{Legendre},
and has been generalised by \gesperrt{Jacobi}. We shall
call it the \gesperrt{Legendre symbol}. Its significance for
number theory will appear presently; first we have
the theorem:

1. \gesperrt{We have} $\left(\dfrac{m}{n}\right)$ \gesperrt{equal to} $+1$ \gesperrt{or equal to} $-1$,
\gesperrt{according as the number of negative ones among the absolutely
least residues of} $m\nu$ \gesperrt{is even or odd.}

For $m = 1$, as we have already seen, the positive sign
holds in formula (2); hence we have, as formula
(4) also shows directly,

\[
2.\quad \left(\frac{1}{n}\right) = +1.
\]

If we change $m$ into $-m$, the sign changes in all factors
of the product on the left-hand side of (4); since
the number of factors is $\dfrac{n-1}{2}$, it follows that

\[
3.\quad \left(\frac{-m}{n}\right) = (-1)^{\frac{n-1}{2}} \left(\frac{m}{n}\right)
\]

\gesperrt{and for} $m = 1$
%CONT>>
